\section{Generated Answers as a Free Compatible Completion}
\label{sec:free-completion}

This section records the algebraic substrate used by the observational results.
Its role is foundational rather than novel: the normalized construction is a
concrete binary presentation of a free compatible completion.

Let $D$ have carrier $A$, binary operation symbols $\Sigma$, and evaluator
\[
  \eval_D:\Sigma\to A\to A\to\operatorname{Option}(A).
\]
Expressions and raw Answers are generated by
\[
 e::=\operatorname{val}(a)\mid\operatorname{app}(f,e_1,e_2),
 \qquad
 r::=\old(a)\mid\susp(f,r_1,r_2).
\]
Here $A$ is consistently called the \emph{base} in the mathematical prose.
The symbol $\old(a)$ is retained only as the constructor name for the embedded
base leaf in the formal Answer syntax, matching the Lean implementation; it
does not introduce a second mathematical notion distinct from the base.
Define the one-step resolution operation by
\[
\operatorname{liftOp}_D(f,x,y)=
\begin{cases}
 \old(c),&x=\old(a),\ y=\old(b),\ \eval_D(f,a,b)=\operatorname{some}(c),\\
 \susp(f,x,y),&\text{otherwise}.
\end{cases}
\]
The evaluator $\Res_D$ recursively resolves both children and then applies
$\operatorname{liftOp}_D$.  The generated Answer algebra $\RAlg{D}$ is the
image of $\Res_D$, with base embedding $\eta_0(a)=\old(a)$ and operations
induced by $\operatorname{liftOp}_D$.

\begin{proposition}[Conservativity on the base]
\label{prop:generated-conservativity}
If $\eval_D(f,a,b)=\operatorname{some}(c)$, then
\[
 f_{\RAlg{D}}(\eta_0(a),\eta_0(b))=\eta_0(c).
\]
The base embedding $\eta_0$ is injective.
\end{proposition}

A compatible total algebra over $D$ consists of a total $\Sigma$-algebra $T$
and a map $i:A\to T$ satisfying
\[
 \eval_D(f,a,b)=\operatorname{some}(c)
 \Longrightarrow f_T(i(a),i(b))=i(c).
\]
The map need not be injective for the free universal property.

Every raw Answer has a canonical fold into $T$:
\[
 \Phi_T(\old(a))=i(a),\qquad
 \Phi_T(\susp(f,x,y))=f_T(\Phi_T(x),\Phi_T(y)).
\]
Compatibility is exactly what is needed to show that this fold respects the
resolution step.

\begin{theorem}[Free compatible completion]\label{thm:free-completion}
For every compatible total algebra $T$ over $D$, there is a unique homomorphism
\[
 \Phi_T:\RAlg{D}\to T
\]
with $\Phi_T\circ\eta_0=i$.
\end{theorem}

The expression algebra modulo equality after resolution is likewise
canonically equivalent to the generated Answer algebra.

\begin{remark}[Positioning]
The theorem above should be read as a concrete normalized presentation of an
established free-completion phenomenon, not as the main novelty claim.  Its
purpose in this paper is to provide canonical generated syntax and canonical
interpretation maps into compatible observers.  All later separation and
completion statements are relative to those interpretation maps.
\end{remark}

Generated Answers are normal forms: a suspended node cannot have two old
children on which the base evaluator is defined, because such a node would have
reduced.  This simple normalization invariant is the key local fact used by the
finite-pattern observers in the next section.
